\documentclass[a4paper,11pt]{article}
\usepackage[a4paper,margin=2cm]{geometry}
\usepackage[T1]{fontenc}
\usepackage[utf8]{inputenc}
\usepackage[ngerman,english]{babel}
\usepackage{lmodern}
\usepackage{microtype}
\usepackage{xcolor}
\usepackage{parskip}
\usepackage{enumitem}
\usepackage[most]{tcolorbox}
\usepackage{titlesec}
\usepackage[hidelinks]{hyperref}
\definecolor{HeaderBlueDark}{RGB}{70,110,170}
\definecolor{RowBlueLight}{RGB}{235,243,255}
\definecolor{RowBlueDark}{RGB}{220,233,250}
\definecolor{SoftGray}{RGB}{90,90,90}
\setlength{\parindent}{0pt}
\setlist[itemize]{leftmargin=1.4em,itemsep=0.35em,topsep=0.35em}
\linespread{1.06}
\titleformat{\section}[block]{\Large\bfseries\color{HeaderBlueDark}}{}{0pt}{}
\titlespacing{\section}{0pt}{1.3em}{0.8em}
\newtcolorbox{exbox}{enhanced,breakable,colback=RowBlueLight,colframe=RowBlueDark,boxrule=0.4pt,arc=1.5mm,left=1.2mm,right=1.2mm,top=1mm,bottom=1mm,title={Beispiele \textnormal{(Examples)}},fonttitle=\bfseries,colbacktitle=RowBlueDark,coltitle=black}
\NewDocumentEnvironment{grammarentry}{mm+b}{%
  \begin{tcolorbox}[enhanced,breakable,colback=white,colframe=HeaderBlueDark,boxrule=0.6pt,arc=1.5mm,left=1.5mm,right=1.5mm,top=1mm,bottom=1mm,colbacktitle=HeaderBlueDark,coltitle=white,fonttitle=\bfseries,title={#1\hfill\normalfont\footnotesize #2}]
  #3
  \end{tcolorbox}
}{}
\newcommand{\GermanText}[1]{\textbf{Deutsch: }#1\par}
\newcommand{\EnglishText}[1]{{\small\color{SoftGray}\textbf{English: }#1\par}}
\newcommand{\ExampleItem}[2]{\item \textbf{#1}\\{\small\color{SoftGray}#2}}
\title{Die Konjunktion \glqq obwohl\grqq{} \textnormal{(The conjunction \glqq obwohl\grqq)}}
\date{}
\begin{document}
\selectlanguage{ngerman}\maketitle\tableofcontents\newpage

\section{Grundbedeutung \textnormal{(Basic meaning)}}
\begin{grammarentry}{Bedeutung von \textbf{obwohl}}{Einräumung / unerwarteter Gegensatz}
\GermanText{\textbf{Obwohl} leitet einen Nebensatz ein und nennt einen Umstand, der normalerweise gegen die Aussage des Hauptsatzes spricht. Der Hauptsatz beschreibt, was trotzdem geschieht oder gilt.}
\EnglishText{\textbf{Obwohl} introduces a subordinate clause and names a circumstance that would normally count against the main-clause statement. The main clause says what happens or is true anyway.}
\GermanText{Übliche Übersetzungen: \textbf{obwohl}, \textbf{auch wenn}, englisch \textbf{although / even though}.}
\EnglishText{Common equivalents: \textbf{although}, \textbf{even though}.}
\end{grammarentry}
\begin{exbox}\begin{itemize}
\ExampleItem{Obwohl es regnet, gehen wir spazieren.}{Although it is raining, we are going for a walk.}
\ExampleItem{Sie ging zur Arbeit, obwohl sie krank war.}{She went to work even though she was ill.}
\end{itemize}\end{exbox}

\section{Wortstellung: Verb am Ende \textnormal{(Word order: verb at the end)}}
\begin{grammarentry}{Subordinierende Konjunktion}{Nebensatz}
\GermanText{\textbf{Obwohl} ist eine subordinierende Konjunktion (Subjunktion). Im eingeleiteten Nebensatz steht das konjugierte Verb am Ende. Bei mehreren Verbteilen steht das finite Verb normalerweise ganz zuletzt.}
\EnglishText{\textbf{Obwohl} is a subordinating conjunction. In the subordinate clause it introduces, the finite verb goes to the end. With multiple verb parts, the finite verb normally comes last.}
\GermanText{Grundmuster: \textbf{obwohl + Subjekt + ... + Verb}.}
\EnglishText{Basic pattern: \textbf{obwohl + subject + ... + verb}.}
\end{grammarentry}
\begin{exbox}\begin{itemize}
\ExampleItem{Obwohl der Bus spät kommt, warten wir.}{Although the bus is late, we are waiting.}
\ExampleItem{Obwohl sie viel gelernt hat, ist sie nervös.}{Although she has studied a lot, she is nervous.}
\ExampleItem{Obwohl er früher hätte gehen können, blieb er bis zum Ende.}{Although he could have left earlier, he stayed until the end.}
\end{itemize}\end{exbox}

\section{Stellung des Nebensatzes \textnormal{(Position of the subordinate clause)}}
\begin{grammarentry}{Nebensatz zuerst}{Verbzweitstellung im Hauptsatz}
\GermanText{Steht der \textbf{obwohl}-Satz zuerst, besetzt er die erste Position des Gesamtsatzes. Nach dem Komma folgt im Hauptsatz direkt das konjugierte Verb.}
\EnglishText{When the \textbf{obwohl}-clause comes first, it occupies the first position in the full sentence. After the comma, the finite verb in the main clause comes next.}
\end{grammarentry}
\begin{exbox}\begin{itemize}
\ExampleItem{Obwohl ich müde bin, arbeite ich weiter.}{Although I am tired, I keep working.}
\ExampleItem{Ich arbeite weiter, obwohl ich müde bin.}{I keep working although I am tired.}
\ExampleItem{Obwohl die Fahrt lang war, kamen wir pünktlich an.}{Although the journey was long, we arrived on time.}
\end{itemize}\end{exbox}

\section{Kommasetzung \textnormal{(Comma use)}}
\begin{grammarentry}{Komma vor und nach dem Nebensatz}{obligatorische Abtrennung}
\GermanText{Ein \textbf{obwohl}-Nebensatz wird durch ein Komma vom Hauptsatz getrennt. Steht der Nebensatz in der Mitte, wird er in der Regel von zwei Kommas eingeschlossen.}
\EnglishText{An \textbf{obwohl}-clause is separated from the main clause by a comma. When it is embedded in the middle, it is normally enclosed by two commas.}
\end{grammarentry}
\begin{exbox}\begin{itemize}
\ExampleItem{Obwohl der Weg weit ist, gehe ich zu Fuß.}{Although the way is long, I am walking.}
\ExampleItem{Ich gehe, obwohl der Weg weit ist, heute zu Fuß.}{I am walking today, although the way is long.}
\end{itemize}\end{exbox}

\section{Indikativ und Konjunktiv \textnormal{(Indicative and subjunctive)}}
\begin{grammarentry}{Tatsache oder angenommener Fall}{Modus}
\GermanText{Für einen als wirklich dargestellten Umstand steht normalerweise der Indikativ. Der Konjunktiv II ist möglich, wenn der Umstand als hypothetisch oder nicht wirklich dargestellt wird.}
\EnglishText{The indicative is normally used for a circumstance presented as real. The subjunctive II is possible when the circumstance is presented as hypothetical or unreal.}
\end{grammarentry}
\begin{exbox}\begin{itemize}
\ExampleItem{Obwohl er krank ist, geht er zur Arbeit.}{Although he is ill, he is going to work.}
\ExampleItem{Obwohl er krank wäre, würde er wahrscheinlich zur Arbeit gehen.}{Even if he were ill, he would probably go to work.}
\end{itemize}\end{exbox}

\section{Unterschied zu \glqq trotz\grqq{} und \glqq trotzdem\grqq{} \textnormal{(Difference from \glqq trotz\grqq{} and \glqq trotzdem\grqq)}}
\begin{grammarentry}{Ähnliche Bedeutung, andere Konstruktion}{Vergleich}
\GermanText{\textbf{Obwohl} verlangt einen vollständigen Nebensatz. \textbf{Trotz} steht vor einer Nominalgruppe, normalerweise im Genitiv. \textbf{Trotzdem} verbindet Aussagen als Konjunktionaladverb; danach steht das finite Verb auf Position 2.}
\EnglishText{\textbf{Obwohl} takes a full subordinate clause. \textbf{Trotz} precedes a noun phrase, normally in the genitive. \textbf{Trotzdem} connects statements as a conjunctive adverb; the finite verb then stays in position two.}
\end{grammarentry}
\begin{exbox}\begin{itemize}
\ExampleItem{Obwohl es regnete, gingen wir spazieren.}{Although it rained, we went for a walk.}
\ExampleItem{Trotz des Regens gingen wir spazieren.}{Despite the rain, we went for a walk.}
\ExampleItem{Es regnete. Trotzdem gingen wir spazieren.}{It rained. Nevertheless, we went for a walk.}
\end{itemize}\end{exbox}

\section{Nahe Ausdrücke \textnormal{(Related expressions)}}
\begin{grammarentry}{\textbf{obgleich}, \textbf{obschon}, \textbf{wenngleich}}{subordinierende Konjunktionen}
\GermanText{\textbf{Obgleich}, \textbf{obschon} und \textbf{wenngleich} haben eine ähnliche Bedeutung und dieselbe Grundwortstellung wie \textbf{obwohl}. Sie klingen oft formeller oder schriftsprachlicher; \textbf{obwohl} ist im Alltag die häufigste Form.}
\EnglishText{\textbf{Obgleich}, \textbf{obschon}, and \textbf{wenngleich} have a similar meaning and the same basic word order as \textbf{obwohl}. They often sound more formal or literary; \textbf{obwohl} is the most common everyday form.}
\end{grammarentry}
\begin{exbox}\begin{itemize}
\ExampleItem{Obgleich die Aufgabe schwierig war, löste er sie schnell.}{Although the task was difficult, he solved it quickly.}
\ExampleItem{Wenngleich die Lösung nicht perfekt ist, hilft sie uns weiter.}{Although the solution is not perfect, it helps us move forward.}
\end{itemize}\end{exbox}

\section{Zusammenfassung \textnormal{(Summary)}}
\begin{grammarentry}{Merksatz}{kurz und wichtig}
\GermanText{\textbf{Obwohl + Nebensatz, Hauptsatz}; das konjugierte Verb steht im obwohl-Satz am Ende. Vor den Satz kommt ein Komma.}
\EnglishText{\textbf{Obwohl + subordinate clause, main clause}; the finite verb goes to the end of the obwohl-clause. Put a comma before the main clause.}
\end{grammarentry}
\end{document}
